\begin{textsample}{POS Dim 1 – vem\_ed\_19 – Score 11.00 – vem\_ed\_19/t091.txt}  \label{ex:f1_pos_051}
\textbf{continuação} Alguns sinais de \textbf{que} um projeto não vai bem: professores \textbf{que} não participam, não detalham o projeto; docentes \textbf{que} culpam o parceiro ou não buscam soluções / bons resultados para o projeto - considerado como atividade pontual sem continuidade e sem visão das relações interinstitucionais; professores \textbf{que} não escutam seus alunos, não \textbf{seguem} o planejado ou mudam os planos no meio do projeto sem consultar parceiros e não buscam um “justo meio-termo”; estudantes \textbf{que} não interagem, confusos ou em constante \textbf{conflito}, sem alguém para coordená-los, ou ainda insensíveis à cultura do outro país, vista sem respeito e com estereótipos.
Do \textbf{ponto} de vista \textbf{institucional}, portanto, \textbf{é} imprescindível \textbf{contar} com um coordenador \textbf{que} ajude a resolver os problemas \textbf{locais}.
Em resumo: para \textbf{expandir} os Intercâmbios Virtuais com qualidade, \textbf{é} preciso ter uma equipe \textbf{que} gerencie o design, a implantação e o \textbf{desenvolvimento} desses projetos.

% matched lemmas: conflito, contar, continuação, desenvolvimento, expandir, institucional, local, ponto, que, seguir, ser
\end{textsample}
